\subsection*{Projection, moments, and ordinary least squares}
Unless stated otherwise in this section, $X\in\R^{n\times p}$ is fixed
and has full column rank, and $Y=X\beta+\varepsilon$ with
$\E\varepsilon=0$ and $\Cov(\varepsilon)=\sigma^2 I_n$, $\sigma^2>0$.
Residual variance estimation requires $n>p$.
\begin{problem}[The residual projector]\label{ex:11:residual-projector}
Let $X\in\R^{n\times p}$ have rank $p$, and set
$H=X(X^\top X)^{-1}X^\top$, $M=I-H$.
Prove that $H,M$ are symmetric idempotent, $HX=X$, $MX=0$, $HM=0$,
and that their ranks and traces are $p$ and $n-p$.
\end{problem}
\begin{solution}
Put $K=X^\top X$. Since $K$ is symmetric positive definite, $K^{-1}$
is symmetric. Hence $H^\top=H$, and
$H^2=XK^{-1}KK^{-1}X^\top=H$. Also $HX=XK^{-1}K=X$.
It follows that $M^\top=M$, $M^2=I-2H+H^2=M$, $MX=0$,
and $HM=H-H^2=0$. The range of $H$ is exactly $\operatorname{col}(X)$:
its outputs lie there and it fixes every column of $X$. Thus its rank
is $p$. The complementary range of $M$ is $\ker X^\top$, of dimension
$n-p$. An idempotent symmetric matrix has eigenvalues $0$ and $1$,
because $Av=\lambda v$ gives $\lambda^2v=\lambda v$.
Its trace and rank both count its eigenvalues $1$, proving the trace claims.
\end{solution}

\begin{problem}[Characterizing a projector by its action]
Let $A=A^\top=A^2$ and let $B\in\R^{n\times m}$ have full column rank.
Prove (a) if $AB=B$ and $\operatorname{rank}A=m$, then
$A=B(B^\top B)^{-1}B^\top$; (b) if $AB=0$ and
$\operatorname{rank}A+m=n$, then $A=I-B(B^\top B)^{-1}B^\top$.
\end{problem}
\begin{solution}
(a) The equality $AB=B$ puts $\operatorname{col}(B)$ inside
$\operatorname{col}(A)$. Their dimensions are both $m$, so the subspaces
coincide. A symmetric idempotent matrix is the orthogonal projector onto
its range: it fixes its range and kills its orthogonal complement.
Both $A$ and $B(B^\top B)^{-1}B^\top$ have this action, hence are equal.
(b) The matrix $I-A$ is symmetric idempotent, fixes $B$, and has rank
$n-\operatorname{rank}A=m$. Apply (a) to $I-A$ and rearrange.
\end{solution}

\begin{problem}[The unique least-squares minimizer]
For full-column-rank $X$ and arbitrary $y$, find the unique minimizer of
$\norm{y-Xb}^2$ and prove global minimality.
\end{problem}
\begin{solution}
Set $b_0=(X^\top X)^{-1}X^\top y$ and $r=y-Xb_0$.
Then $X^\top r=0$. For any $b$,
\[
\norm{y-Xb}^2=\norm{r-X(b-b_0)}^2
=\norm r^2+(b-b_0)^\top X^\top X(b-b_0).
\]
The second term is positive unless $b=b_0$, since $X$ has trivial kernel.
Thus $b_0$ is the unique global minimizer. This proof checks global
minimality directly, without a second-derivative assumption.
\end{solution}

\begin{problem}[Unbiased coefficients and covariance]
Under $Y=X\beta+\varepsilon$, $\E\varepsilon=0$ and
$\Cov(\varepsilon)=\sigma^2I$, prove that the least-squares coefficient
estimator is unbiased and find its covariance.
\end{problem}
\begin{solution}
Let $L=(X^\top X)^{-1}X^\top$. Since $LX=I$,
$\widehat\beta=\beta+L\varepsilon$. Thus
$\E\widehat\beta=\beta+L\E\varepsilon=\beta$, and
\[
\Cov(\widehat\beta)=L(\sigma^2I)L^\top
=\sigma^2(X^\top X)^{-1}X^\top X(X^\top X)^{-1}
=\sigma^2(X^\top X)^{-1}.
\]
Full column rank is needed even for defining this formula.
\end{solution}

\begin{problem}[Unbiased residual variance]
For $p<n$, prove that $s^2=Y^\top MY/(n-p)$ is unbiased for $\sigma^2$.
\end{problem}
\begin{solution}
Since $MX=0$, $MY=M\varepsilon$, and
$Y^\top MY=\varepsilon^\top M\varepsilon$.
Expanding by entries gives
$\E(\varepsilon^\top M\varepsilon)=\sum_{i,j}m_{ij}\E(\varepsilon_i\varepsilon_j)
=\sigma^2\sum_i m_{ii}=\sigma^2(n-p)$.
Division proves the result. If $p=n$, then $M=0$, so SSE is always zero
and this estimator cannot be defined by dividing by residual degrees
of freedom.
\end{solution}

\begin{problem}[Leverages and residual variances]
Prove $0\leqslant h_{ii}\leqslant1$, $\sum_i h_{ii}=p$, and
$\Var(e_i)=\sigma^2(1-h_{ii})$ for $e=MY$.
\end{problem}
\begin{solution}
For the coordinate column $d_i$, $h_{ii}=\norm{Hd_i}^2\geqslant0$.
The orthogonal sum $d_i=Hd_i+Md_i$ gives
$1=h_{ii}+\norm{Md_i}^2$, proving the upper bound.
The sum equals $\tr H=p$. Finally $e=M\varepsilon$ and
$\Cov(e)=\sigma^2MM^\top=\sigma^2M$, whose $i$th diagonal is the claimed
variance. No independence of the residual coordinates is implied.
\end{solution}

\begin{problem}[A numerical fit and its covariance]\label{ex:11:numerical}
For $X=(\mathbf1\ t)$, $t=(-3,-1,1,3)^\top$, and $y=(1,2,2,5)^\top$,
compute the least-squares coefficients, fitted values, residuals, SSE,
$s^2$, and the estimated covariance of the coefficient estimator.
Verify the two residual orthogonality constraints.
\end{problem}
\begin{solution}
$X^\top X=\operatorname{diag}(4,20)$ and $X^\top y=(10,12)^\top$, so
$\widehat\beta=(5/2,3/5)^\top$.
The fit is $(7,19,31,43)^\top/10$ and the residual is
$(3,1,-11,7)^\top/10$. Its entries sum to zero, and its dot product
with $t$ is $(-9-1-11+21)/10=0$.
The SSE is $(9+1+121+49)/100=9/5$. With $n-p=2$,
$s^2=9/10$ and the estimated covariance is
$s^2(X^\top X)^{-1}=\operatorname{diag}(9/40,9/200)$.
\end{solution}

\begin{problem}[Gauss--Markov in covariance order]
State and prove the covariance-order optimality of ordinary least squares
among linear unbiased coefficient estimators, including the equality case.
\end{problem}
\begin{solution}
For $B$ with $BX=I$, put $D=B-L$, where $L=(X^\top X)^{-1}X^\top$.
Then $DX=0$, and $LD^\top=(X^\top X)^{-1}(DX)^\top=0$.
Consequently
$\Cov(BY)=\sigma^2BB^\top=\sigma^2LL^\top+\sigma^2DD^\top$.
For any $c$, the excess variance of $c^\top BY$ is
$\sigma^2\norm{D^\top c}^2\geqslant0$, proving the matrix order.
Equality of the complete covariance matrices gives
$0=\tr(DD^\top)=\sum_{i,j}d_{ij}^2$, so $B=L$; the converse is immediate.
The hypotheses are fixed full-column-rank $X$, zero mean errors, and
covariance $\sigma^2I$ with $\sigma^2>0$.
\end{solution}

\begin{problem}[Minimum total variance as a matrix problem]
For a given full-column-rank $X\in\R^{n\times p}$, find the
$p\times n$ matrix $G$ minimizing $\tr(GG^\top)$ subject to $GX=I_p$.
\end{problem}
\begin{solution}
Let $L=(X^\top X)^{-1}X^\top$. Every feasible $G$ is $L+D$ with $DX=0$.
The cross products $LD^\top$ and $DL^\top$ vanish, so
\[
\tr(GG^\top)=\tr(LL^\top)+\tr(DD^\top)
=\tr((X^\top X)^{-1})+\sum_{i,j}d_{ij}^2.
\]
The second term is nonnegative and is zero exactly when $D=0$.
Thus the unique minimizer is $G=L$. This solves the stated scalar trace
minimization; the preceding covariance-order theorem is stronger.
\end{solution}

\begin{problem}[Centered sum-of-squares decomposition]
When an intercept is present, prove $\mathrm{SST}=\mathrm{SSR}+\mathrm{SSE}$,
with all three sums centered or uncentered as defined in the notes.
\end{problem}
\begin{solution}
Put $J=\mathbf1\mathbf1^\top/n$. The intercept gives $HJ=JH=J$.
Then $H-J$ and $M=I-H$ are orthogonal projectors with
$(H-J)M=0$ and $(H-J)+M=I-J$. Hence
$(I-J)y=(H-J)y+My$ is an orthogonal sum. Taking squared lengths gives
\[
\norm{y-\bar y\mathbf1}^2
=\norm{Hy-\bar y\mathbf1}^2+\norm{My}^2.
\]
These are exactly SST, SSR and SSE. The argument depends on the intercept;
without it, $HJ=J$ need not hold.
\end{solution}

\begin{problem}[Nested model spaces]
For nested model spaces $\mathcal S_0\subseteq\mathcal S_1$, prove that
their residual sums of squares satisfy $\mathrm{SSE}_0\geqslant\mathrm{SSE}_1$.
\end{problem}
\begin{solution}
Let $H_i$ be the respective orthogonal projectors. Since $H_1$ fixes
$\mathcal S_0$, $H_1H_0=H_0$; transposition gives $H_0H_1=H_0$.
Thus $H_1-H_0$ is a symmetric idempotent matrix orthogonal to $I-H_1$.
The residual decomposition
$(I-H_0)y=(H_1-H_0)y+(I-H_1)y$ gives
$\mathrm{SSE}_0-\mathrm{SSE}_1=\norm{(H_1-H_0)y}^2\geqslant0$.
Equality holds precisely when the added model directions receive no
component of the observed response.
\end{solution}

\begin{problem}[Partitioned regression]\label{ex:11:fwl}
Let $X=(X_1\ X_2)$ have full column rank and
$M_2=I-X_2(X_2^\top X_2)^{-1}X_2^\top$.
Prove that the first coefficient block is
$(X_1^\top M_2X_1)^{-1}X_1^\top M_2Y$, and that its residualized
regression has the same residuals as the full regression.
\end{problem}
\begin{solution}
If $M_2X_1v=0$, then $X_1v=X_2u$ for some $u$. Full column rank of
$(X_1\ X_2)$ forces $v=u=0$, so $X_1^\top M_2X_1$ is positive definite.
For fixed $b_1$, the least-squares $b_2$ equals
$(X_2^\top X_2)^{-1}X_2^\top(Y-X_1b_1)$.
Substitution makes the residual $M_2Y-M_2X_1b_1$.
The normal equations for this remaining full-column-rank problem are
$X_1^\top M_2X_1b_1=X_1^\top M_2Y$, giving the formula.
At its solution the displayed residual is exactly the full residual,
because the eliminated $b_2$ was its unique optimal value for that $b_1$.
\end{solution}

\begin{problem}[Omitted-variable bias]
Suppose $\E Y=X_1\beta_1+X_2\beta_2$, but fit only the full-column-rank
design $X_1$. Find the expectation of the reduced coefficient estimator
and the condition for it to be unbiased for $\beta_1$.
\end{problem}
\begin{solution}
The reduced estimator is $\widetilde\gamma=(X_1^\top X_1)^{-1}X_1^\top Y$.
Taking its expectation gives
\[
\E\widetilde\gamma=\beta_1+
(X_1^\top X_1)^{-1}X_1^\top X_2\beta_2.
\]
For a specified $\beta_2$, unbiasedness holds exactly when
$X_1^\top X_2\beta_2=0$. For all possible $\beta_2$, it holds exactly
when $X_1^\top X_2=0$. A zero omitted coefficient vector is sufficient
in the first statement; it is not necessary, since cancellations in
the matrix-vector product are possible.
\end{solution}

\begin{problem}[Prediction variance]
Let $Y_0=x_0^\top\beta+\varepsilon_0$ be a new observation with
$\E\varepsilon_0=0$, $\Var(\varepsilon_0)=\sigma^2$, and
$\Cov(\varepsilon_0,\varepsilon)=0$. For
$\widehat Y_0=x_0^\top\widehat\beta$, compare the variance of estimating
the mean $x_0^\top\beta$ with that of predicting $Y_0$.
\end{problem}
\begin{solution}
The first error is $x_0^\top(\widehat\beta-\beta)$, with variance
$\sigma^2x_0^\top(X^\top X)^{-1}x_0$.
The prediction error subtracts $\varepsilon_0$. Its cross-covariance
with the coefficient error is zero by the assumed error cross-covariance,
so its variance is
$\sigma^2[1+x_0^\top(X^\top X)^{-1}x_0]$.
It exceeds the mean-estimation variance by $\sigma^2>0$, the variability
of the new observation about its own mean.
\end{solution}

\begin{problem}[Adjusted explained variation]
For two full-rank regressions sharing an intercept and differing by one
added column, assume SST is positive and $n-p-1>0$, where $p$ is the
reduced parameter count. Prove that $R^2$ cannot decrease. Define
\[
\overline R^2=1-\frac{\mathrm{SSE}/(n-\text{parameter count})}{\mathrm{SST}/(n-1)}.
\]
Show that $\overline R^2$ strictly increases exactly when the added
coefficient's $|t|$ exceeds $1$, where $t$ uses the full-model residual
variance and is defined when that variance is positive.
\end{problem}
\begin{solution}
Nested-model geometry gives $\mathrm{SSE}_F\leqslant\mathrm{SSE}_R$.
Since SST is the same positive number, $R^2=1-\mathrm{SSE}/\mathrm{SST}$
cannot decrease. Let $z$ be the new column residualized against the old
design; full rank gives $z\ne0$. Its fitted coefficient is
$\widehat\gamma=z^\top y/(z^\top z)$, and the SSE decrease is
$\Delta=\widehat\gamma^2z^\top z$.
Put $d=n-p-1$ and $s_F^2=\mathrm{SSE}_F/d>0$.
Then $t^2=\widehat\gamma^2/(s_F^2/(z^\top z))=\Delta/s_F^2$.
The adjusted value increases exactly when
$\mathrm{SSE}_F/d<\mathrm{SSE}_R/(d+1)$, or
$\mathrm{SSE}_F<d\Delta$, equivalently $t^2>1$.
This is an algebraic identity; no distributional interpretation of $t$
is needed. If the full SSE is zero, $t$ in this form is undefined.
\end{solution}

\subsection*{Weighted and constrained estimation}
\begin{problem}[Generalized least squares]\label{ex:11:gls}
For full-column-rank $X$ and positive definite $\Omega$, find the unique
minimizer of $(y-Xb)^\top\Omega^{-1}(y-Xb)$.
\end{problem}
\begin{solution}
Set $Z=\Omega^{-1/2}X$, $v=\Omega^{-1/2}y$. Since the whitening matrix
is invertible, $Z$ has full column rank. The objective becomes
$\norm{v-Zb}^2$, whose unique minimizer is
\[
b_0=(Z^\top Z)^{-1}Z^\top v
=(X^\top\Omega^{-1}X)^{-1}X^\top\Omega^{-1}y.
\]
Equivalently, for every displacement $d$, the objective increase from
$b_0$ to $b_0+d$ is $d^\top X^\top\Omega^{-1}Xd$, strictly positive
for $d\ne0$. This also verifies global uniqueness.
\end{solution}

\begin{problem}[Actual OLS covariance and optimal GLS]
Under $\Cov(\varepsilon)=\sigma^2\Omega\succ0$, prove that OLS is still
unbiased, derive its actual covariance, and prove that GLS is best linear
unbiased in covariance order.
\end{problem}
\begin{solution}
With $K=X^\top X$, $L=K^{-1}X^\top$ still satisfies $LX=I$.
Thus $\E(LY)=\beta$ and
$\Cov(LY)=\sigma^2K^{-1}X^\top\Omega XK^{-1}$.
Let $L_\Omega=(X^\top\Omega^{-1}X)^{-1}X^\top\Omega^{-1}$.
It satisfies $L_\Omega X=I$ and has covariance
$\sigma^2(X^\top\Omega^{-1}X)^{-1}$.
For any unbiased $BY$, put $D=B-L_\Omega$; then $DX=0$ and
$L_\Omega\Omega D^\top=0$. Expansion gives
$\Cov(BY)-\Cov(L_\Omega Y)=\sigma^2D\Omega D^\top\succeq0$.
The sandwich formula is the correct OLS covariance under this model;
the spherical formula generally fails.
\end{solution}

\begin{problem}[Minimum total variance for a general target]
Let $X\in\R^{n\times p}$ have full column rank, $\Omega\succ0$, and
$W\in\R^{m\times p}$ be fixed. Minimize $\tr(G\Omega G^\top)$ over
$G\in\R^{m\times n}$ subject to $GX=W$.
\end{problem}
\begin{solution}
Let $K=X^\top\Omega^{-1}X$ and $G_0=WK^{-1}X^\top\Omega^{-1}$.
It is feasible because $G_0X=W$. Every feasible $G$ is $G_0+D$ with
$DX=0$, and $G_0\Omega D^\top=WK^{-1}(DX)^\top=0$.
Therefore
$\tr(G\Omega G^\top)=\tr(WK^{-1}W^\top)+\tr(D\Omega D^\top)$.
The last term is $\norm{D\Omega^{1/2}}_F^2$, zero only when $D=0$.
Thus $G_0$ is the unique minimizer, including when $W$ has dependent or
zero rows.
\end{solution}

\begin{problem}[A matrix Cauchy--Schwarz inequality]
Suppose $A,B\in\R^{n\times p}$ and $A^\top B=I_p$. Prove
$A^\top A\succeq(B^\top B)^{-1}$ and
$B^\top B\succeq(A^\top A)^{-1}$, and characterize equality.
\end{problem}
\begin{solution}
$Bx=0$ implies $x=A^\top Bx=0$, so $B^\top B$ is positive definite.
Let $P=B(B^\top B)^{-1}B^\top$. Since $B^\top A=I$,
\[
A^\top A-(B^\top B)^{-1}=A^\top(I-P)A
=[(I-P)A]^\top[(I-P)A]\succeq0.
\]
Equality holds exactly when $(I-P)A=0$, or
$A=PA=B(B^\top B)^{-1}$.
For the second inequality, $B^\top A=I$ also gives full column rank of
$A$. With $Q=A(A^\top A)^{-1}A^\top$, the difference is
$B^\top(I-Q)B\succeq0$, zero exactly when $B=A(A^\top A)^{-1}$.
The equality conditions are equivalent by substitution into either
Gram matrix.
\end{solution}

\begin{problem}[Two proofs of the OLS--GLS inequality]
Prove, for full-column-rank $X$ and $V\succ0$,
\[
(X^\top X)^{-1}X^\top VX(X^\top X)^{-1}
\succeq(X^\top V^{-1}X)^{-1},
\]
first using the preceding matrix Cauchy--Schwarz inequality, then using
the projector onto $\operatorname{col}(V^{-1/2}X)$. Explain its statistical meaning.
\end{problem}
\begin{solution}
Set $A=V^{1/2}X(X^\top X)^{-1}$ and $B=V^{-1/2}X$. Then $A^\top B=I$,
and the preceding inequality gives exactly the displayed formula.
For a second proof put
$P=V^{-1/2}X(X^\top V^{-1}X)^{-1}X^\top V^{-1/2}$.
It is an orthogonal projector, so
\[
V^{1/2}(I-P)V^{1/2}
=V-X(X^\top V^{-1}X)^{-1}X^\top\succeq0.
\]
Multiplying on both sides by the transpose pair
$(X^\top X)^{-1}X^\top$ and $X(X^\top X)^{-1}$ gives the claim.
When $\Cov(Y)=V$, the two sides are the actual covariances of OLS and
GLS. Thus GLS never has a larger variance for any coefficient contrast.
Equality can occur; OLS need not be strictly worse.
\end{solution}

\begin{problem}[$*$ When OLS and GLS agree]\label{ex:11:agreement}
Let $K=X^\top X$, $J=X^\top V^{-1}X$, $M=I-XK^{-1}X^\top$, with
$X$ of full column rank and $V\succ0$. Prove the equivalence of
\[
X^\top V^{-1}M=0,\qquad
J^{-1}X^\top V^{-1}=K^{-1}X^\top,\qquad
J^{-1}=K^{-1}X^\top VXK^{-1}.
\]
\end{problem}
\begin{solution}
Expanding the first equality gives
$X^\top V^{-1}=JK^{-1}X^\top$; multiplying by $J^{-1}$ proves its
equivalence with the second. Let $L=K^{-1}X^\top$, $L_V=J^{-1}X^\top V^{-1}$.
If $L=L_V$, their covariance matrices $LVL^\top$ and
$L_VVL_V^\top=J^{-1}$ agree, giving the third equality.
Conversely set $D=L-L_V$. Since $DX=0$, both cross terms in
$(L_V+D)V(L_V+D)^\top$ vanish. Thus
\[
K^{-1}X^\top VXK^{-1}-J^{-1}=DVD^\top.
\]
The third equality makes the right side zero. Taking its trace gives
$\norm{DV^{1/2}}_F^2=0$; invertibility of $V^{1/2}$ implies $D=0$,
which is the second equality. This proves all directions, including the
nontrivial implication from equal covariance to equal estimator matrices.
\end{solution}

\begin{problem}[A stronger covariance bound]
For full-column-rank $X$ and $V\succ0$, prove
$X(X^\top V^{-1}X)^{-1}X^\top\preceq V$.
Deduce $(X^\top V^{-1}X)^{-1}\preceq X^\top VX$ when $X^\top X=I$.
\end{problem}
\begin{solution}
The projector $P$ onto $\operatorname{col}(V^{-1/2}X)$ satisfies
$0\preceq P\preceq I$. Multiplication by $V^{1/2}$ on both sides gives
$X(X^\top V^{-1}X)^{-1}X^\top\preceq V$.
If $X^\top X=I$, multiply this inequality on the left by $X^\top$ and
on the right by $X$. The left side reduces to the requested inverse.
Congruence preserves the semidefinite order because its quadratic forms
are the original quadratic forms evaluated on transformed vectors.
\end{solution}

\begin{problem}[Weighted averaging]
For $X=(1,1)^\top$, $y=(1,3)^\top$, and $\Omega=\operatorname{diag}(1,4)$,
compute OLS and GLS, their variances under covariance $\sigma^2\Omega$,
and the GLS fitted matrix and residual. Verify weighted orthogonality.
\end{problem}
\begin{solution}
OLS is $(1+3)/2=2$ and its variance is
$\sigma^2(1/2,1/2)\Omega(1/2,1/2)^\top=5\sigma^2/4$.
For GLS, $X^\top\Omega^{-1}X=5/4$ and
$X^\top\Omega^{-1}y=7/4$, so the estimate is $7/5$ and its variance
is $4\sigma^2/5$.
The fitted matrix is $H_\Omega=\begin{pmatrix}4/5&1/5\\4/5&1/5\end{pmatrix}$,
the fit is $(7/5,7/5)^\top$, and the residual is $(-2/5,8/5)^\top$.
The weighted residual inner product with $X$ is
$-2/5+(1/4)(8/5)=0$. Although $H_\Omega$ is not symmetric, its square
equals itself because its two row entries sum to one.
\end{solution}

\begin{problem}[Constrained weighted least squares]\label{ex:11:constraint}
For full-column-rank $X$, $\Omega\succ0$, and full-row-rank
$R\in\R^{q\times p}$, minimize
$(y-Xb)^\top\Omega^{-1}(y-Xb)$ subject to $Rb=r$.
Find the exact increase above the unconstrained minimum.
\end{problem}
\begin{solution}
Put $K=X^\top\Omega^{-1}X$, $b_0=K^{-1}X^\top\Omega^{-1}y$,
$C=RK^{-1}R^\top$, and $d=Rb_0-r$. Both $K,C$ are positive definite.
The candidate is $b_R=b_0-K^{-1}R^\top C^{-1}d$, and
$Rb_R=Rb_0-CC^{-1}d=r$.
Every feasible $b$ has form $b_R+v$ with $Rv=0$. Weighted normal
equations give $Q(b)=Q(b_0)+(b-b_0)^\top K(b-b_0)$.
Since $v^\top K(b_R-b_0)=-v^\top R^\top C^{-1}d=0$,
\[
Q(b)=Q(b_0)+d^\top C^{-1}d+v^\top Kv.
\]
The last term is strictly positive for $v\ne0$, proving uniqueness,
and the middle term is the exact increase.
\end{solution}

\begin{problem}[Constrained OLS by multipliers]
Derive the constrained ordinary least-squares formula under $Rb=r$
using Lagrange multipliers, where $X$ has full column rank and $R$ has
full row rank. Prove that its SSE is at least the unconstrained SSE.
\end{problem}
\begin{solution}
Use $\mathcal L(b,\lambda)=\frac12\norm{y-Xb}^2+\lambda^\top(Rb-r)$.
Differentiation gives $Kb-X^\top y+R^\top\lambda=0$ and $Rb=r$,
where $K=X^\top X$. Thus $b=b_0-K^{-1}R^\top\lambda$ and
$RK^{-1}R^\top\lambda=Rb_0-r$. Since the coefficient matrix is positive
definite, this determines $\lambda$ uniquely and gives
$b_R=b_0-K^{-1}R^\top(RK^{-1}R^\top)^{-1}(Rb_0-r)$.
To verify that the stationary point is a global constrained minimum,
write any other feasible $b=b_R+v$ with $Rv=0$ and use the orthogonal
completion in the preceding solution. The SSE increase is the
nonnegative quadratic form
$(Rb_0-r)^\top(RK^{-1}R^\top)^{-1}(Rb_0-r)$.
\end{solution}

\begin{problem}[A numerical restriction]
For the data in Exercise~\ref{ex:11:numerical}, impose the restriction
that the slope is zero. Compute the constrained coefficients and SSE,
and verify the increase formula.
\end{problem}
\begin{solution}
The fitted constant is the mean $5/2$. Thus the constrained coefficient
column is $(5/2,0)^\top$ and its residual is
$(-3/2,-1/2,-1/2,5/2)^\top$, with squared length $9$.
Here $R=(0\ 1)$, $K^{-1}=\operatorname{diag}(1/4,1/20)$,
$RK^{-1}R^\top=1/20$ and $Rb_0=3/5$.
The increase formula gives $(3/5)^2/(1/20)=36/5$; adding the unconstrained
SSE $9/5$ gives $9$, as computed directly.
\end{solution}

\begin{problem}[$*$ Rank-deficient estimation of the mean]
Let $X\in\R^{n\times p}$ have arbitrary rank and $\Omega\succ0$.
Minimize $\tr(G\Omega G^\top)$ over $n\times n$ matrices satisfying
$GX=X$.
\end{problem}
\begin{solution}
Put $Z=\Omega^{-1/2}X$, $P=ZZ^+$ and
$G_0=\Omega^{1/2}P\Omega^{-1/2}$.
Then $G_0X=\Omega^{1/2}PZ=X$, so it is feasible.
Every feasible $G$ is $G_0+D$ with $DX=0$.
Let $E=D\Omega^{1/2}$. Then $EZ=0$, which implies $EP=EZZ^+=0$.
The Frobenius cross term between
$G_0\Omega^{1/2}=\Omega^{1/2}P$ and $E$ is
$\tr(\Omega^{1/2}PE^\top)=0$. Therefore
\[
\tr(G\Omega G^\top)=\tr(G_0\Omega G_0^\top)+\norm E_F^2,
\]
uniquely minimized by $D=0$. Using the SVD identity
$Z(Z^\top Z)^+Z^\top=P$, the answer also equals
$G_0=X(X^\top\Omega^{-1}X)^+X^\top\Omega^{-1}$.
If $X=0$, the constraint is vacuous and this formula gives $G_0=0$,
the correct unique minimizer.
\end{solution}

\begin{problem}[$*$ A residual quadratic minimization]\label{ex:11:quadratic-min}
Let $X$ have rank $r$. Minimize $n^{-1}\tr(V^2)$ over symmetric positive
semidefinite $V\in\R^{n\times n}$ satisfying $VX=0$ and $\tr V=1$.
Treat the case $r=n$ as well as $r<n$.
\end{problem}
\begin{solution}
If $r=n$, $X$ spans $\R^n$, so $VX=0$ forces $V=0$, contradicting
$\tr V=1$. Thus the feasible set is empty.
If $r<n$, let $Q$ have orthonormal columns spanning $\ker X^\top$,
of dimension $d=n-r$. Symmetry and $VX=0$ imply
$V=QWQ^\top$ with $W=Q^\top VQ\succeq0$.
Then $\tr W=1$ and $\tr(V^2)=\tr(W^2)$.
If $\lambda_1,\ldots,\lambda_d\geqslant0$ are the eigenvalues of $W$,
Cauchy--Schwarz gives $1=(\sum_i\lambda_i)^2\leqslant d\sum_i\lambda_i^2$.
Equality holds exactly when all $\lambda_i=1/d$.
Hence the unique minimizer is
$V=QQ^\top/d=(I-XX^+)/(n-r)$ and the minimum is $1/[n(n-r)]$.
\end{solution}

\begin{problem}[Complementary spaces and a determinant identity]
Suppose $X\in\R^{n\times p}$ and $A\in\R^{n\times(n-p)}$ both have
full column rank, $A^\top X=0$, and $\Omega\succ0$. Prove
\[
\det(A^\top\Omega A)\det(X^\top X)
=\det(A^\top A)\det(X^\top\Omega^{-1}X)\det\Omega.
\]
\end{problem}
\begin{solution}
Choose orthonormal bases $U,V$ of the two complementary column spaces,
and write $X=UR$, $A=VS$ with square invertible $R,S$.
All factors $\det(R)^2$ and $\det(S)^2$ cancel, reducing the assertion to
$\det(V^\top\Omega V)=\det(U^\top\Omega^{-1}U)\det\Omega$.
In the orthogonal basis $(U\ V)$ write
$\Omega=\begin{pmatrix}B&C\\C^\top&D\end{pmatrix}$.
Here $D\succ0$. Block elimination gives
$\det\Omega=\det D\det(B-CD^{-1}C^\top)$.
Solving the block system for the first block column of the inverse gives
$U^\top\Omega^{-1}U=(B-CD^{-1}C^\top)^{-1}$.
Taking determinants proves the reduced identity and hence the original.
The endpoint cases use the convention that an empty determinant is $1$.
\end{solution}

\subsection*{Random-variable spaces and sensitivity: additional reading}
\begin{problem}[$*$ Square-integrable random variables]
On a fixed probability space, show that the real random variables with
finite second moment form a vector space, and that the finite linear
combinations of any such $Z_1,\ldots,Z_k$ form a subspace.
Identify random variables equal with probability one, and explain the
zero vector in this convention.
\end{problem}
\begin{solution}
The pointwise inequality $(U+V)^2\leqslant2U^2+2V^2$ gives
$\E(U+V)^2\leqslant2\E U^2+2\E V^2<\infty$.
Also $\E(cU)^2=c^2\E U^2<\infty$. Thus addition and scaling preserve
finite second moments. The ordinary pointwise real-number identities
give all vector-space axioms. If representatives are changed on sets
of probability zero, sums and scalar multiples change only on unions
of such sets, again of probability zero, so the operations on classes
are well-defined. The zero class consists of variables zero with
probability one. Finally, linear combinations of the finite list remain
such combinations after adding or scaling their coefficients, so they
form a subspace. Multiple coefficient representations do not make the
random variable itself ambiguous.
\end{solution}

For the following exercises, the differential is the linear part of a
change: $df(x)[h]=\left.\frac{d}{dt}f(x+th)\right|_{t=0}$.
For scalar $f$, its derivative is the row $Df(x)$ satisfying
$df(x)[h]=Df(x)h$. Matrix derivatives may be specified as linear maps
of a perturbation matrix. Product and transpose rules follow entry by
entry from ordinary differentiation. Differentiating $A^{-1}A=I$ gives
$d(A^{-1})=-A^{-1}(dA)A^{-1}$ wherever $A$ is invertible.

\begin{problem}[Linear, quadratic, and bilinear differentiation]
For constant $a,A$, find the differential and derivative of $a^\top x$,
$x^\top Ax$, and $x_1^\top Ax_2$ with respect to the stacked variable
$(x_1^\top,x_2^\top)^\top$.
\end{problem}
\begin{solution}
The first differential is $a^\top dx$, so the derivative is $a^\top$.
For the second,
$d(x^\top Ax)=(dx)^\top Ax+x^\top A\,dx=x^\top(A+A^\top)\,dx$,
so its derivative is $x^\top(A+A^\top)$.
For the third,
$d(x_1^\top Ax_2)=x_2^\top A^\top dx_1+x_1^\top A\,dx_2$,
so the derivative is the block row $(x_2^\top A^\top\ \ x_1^\top A)$.
These expressions include every first-order term; the product of two
perturbations contributes only a second-order term.
\end{solution}

\begin{problem}[Least squares by differentiation]
Derive the ordinary least-squares formula by differentiating
$Q(b)=\norm{y-Xb}^2$, and verify the second-order condition.
\end{problem}
\begin{solution}
Expansion gives $Q(b)=y^\top y-2b^\top X^\top y+b^\top Kb$ with
$K=X^\top X\succ0$.
Thus $dQ=2(Kb-X^\top y)^\top db$, giving the stationary point
$b_0=K^{-1}X^\top y$. The Hessian is $2K\succ0$.
For global rather than merely local minimality,
$Q(b_0+d)-Q(b_0)=d^\top Kd>0$ for every $d\ne0$.
\end{solution}

\begin{problem}[Sensitivity of the coefficient estimator]
With $y$ fixed and $X$ of full column rank, obtain the differential of
$b=(X^\top X)^{-1}X^\top y$ as a linear map of $dX$.
Then give its derivative when only column $x_j$ of $X$ varies.
\end{problem}
\begin{solution}
Let $K=X^\top X$ and $e=y-Xb$.
Differentiate the identity $Kb=X^\top y$:
\[
K\,db+(dX)^\top Xb+X^\top(dX)b=(dX)^\top y.
\]
Solving gives
$db=K^{-1}(dX)^\top e-K^{-1}X^\top(dX)b$.
This is the derivative as a linear map from matrices to columns.
If only column $j$ changes, $dX=(dx_j)d_j^\top$, with $d_j$ the
$j$th coordinate column of $\R^p$. Hence
\[
db=K^{-1}(d_je^\top-b_jX^\top)\,dx_j,
\]
and the displayed $p\times n$ matrix is the derivative with respect
to $x_j^\top$. Full column rank persists in a sufficiently small
neighborhood because $\det(X^\top X)>0$ and is continuous.
\end{solution}

\begin{problem}[Sensitivity of the residual]
For the same problem, obtain the differential of $e=y-Xb$ as a linear
map of $dX$, and its derivative when only $x_j$ varies.
\end{problem}
\begin{solution}
Using $de=-(dX)b-X\,db$ and the previous coefficient differential gives
\[
de=-M(dX)b-XK^{-1}(dX)^\top e.
\]
For $dX=(dx_j)d_j^\top$, collect the terms to get
$de=-(b_jM+XK^{-1}d_je^\top)\,dx_j$.
Thus the derivative is $-(b_jM+XK^{-1}d_je^\top)$.
The first term changes the observed residualized fitted column; the
second accounts for the adjustment of the estimated coefficients.
\end{solution}

\begin{problem}[Sensitivity of generalized least squares]
Let the positive definite matrix $\Omega(t)$ vary differentiably with a
real parameter, while $X,y$ stay fixed. Find the derivative of
$b(t)=(X^\top\Omega^{-1}X)^{-1}X^\top\Omega^{-1}y$.
\end{problem}
\begin{solution}
Put $K=X^\top\Omega^{-1}X$ and $e=y-Xb$.
Differentiating the weighted normal equations
$X^\top\Omega^{-1}(y-Xb)=0$ gives
$K\dot b=X^\top\frac{d\Omega^{-1}}{dt}e$.
The inverse rule gives $\frac{d\Omega^{-1}}{dt}=-\Omega^{-1}\dot\Omega\Omega^{-1}$.
Therefore
\[
\dot b=-K^{-1}X^\top\Omega^{-1}\dot\Omega\Omega^{-1}e.
\]
This answers the scalar-parameter sensitivity question directly; no
vectorization or Kronecker notation is needed.
\end{solution}

\begin{problem}[Sensitivity of the residual projector]
For full-column-rank $X$ and $M=I-X(X^\top X)^{-1}X^\top$, obtain the
differential of $M$ as a linear map of $dX$. Specialize it to perturbing
only column $j$.
\end{problem}
\begin{solution}
With $K=X^\top X$, product and inverse differentiation give
\[
\begin{split}
dM={}&-(dX)K^{-1}X^\top-XK^{-1}(dX)^\top\\
&+XK^{-1}[(dX)^\top X+X^\top(dX)]K^{-1}X^\top.
\end{split}
\]
Collect the first and fourth terms, and the second and third terms:
\[
dM=-M(dX)K^{-1}X^\top-XK^{-1}(dX)^\top M.
\]
If $dX=(dx_j)d_j^\top$, put $v=XK^{-1}d_j$ to obtain
$dM=-M(dx_j)v^\top-v(dx_j)^\top M$.
This explicitly gives the linear derivative for either type of
perturbation and shows that $dM$ is symmetric.
\end{solution}
